\section{Related Work}
\label{sec:related_work}

\subsection{Music Conditioning and Dance Generation}
Music-conditioned motion depends on representations of temporal structure
and musical character. Explicit acoustic descriptors extracted with
librosa~\cite{mcfee2015librosa} provide rhythmic and spectral information,
while pretrained models such as MERT~\cite{li2023mert} and
MuQ~\cite{zhu2025muq} provide learned audio representations.
Encoding observed music and forecasting music that has not yet arrived are
distinct conditioning operations; streaming generation must specify which
audio is available at each decision.

AIST++/FACT~\cite{li2021aist} and Bailando~\cite{siyao2022bailando}
develop music-conditioned choreography through sequence modeling and discrete
motion representations. EDGE~\cite{tseng2023edge} and
LODGE~\cite{li2024lodge} develop diffusion-based generation and long-dance
modeling, while FineDance~\cite{li2023finedance} supports fine-grained full-body
choreography. MEGADance~\cite{yang2025megadance} explores genre-aware generation,
and SoulNet~\cite{li2025soulnet} combines hierarchical motion modeling with
learned music--motion alignment. DiscoForcing~\cite{ji2026discoforcing} extends
audio-driven generation to streaming character control and humanoid deployment.
These developments motivate studying how musical conditions guide a robot-native
generator that continues coherently from committed movement.

Generative music models offer a prospective source of conditioning.
Jukebox~\cite{dhariwal2020jukebox} models continuation in a learned audio space,
and Live Music Models~\cite{gdmlyria2025live} studies continuously generated music
with synchronized control. We use a frozen causal music model,
Magenta RealTime~2~\cite{magenta2026mrt2}, to predict upcoming music states from
received audio. These states supply inferred musical lookahead, not access to
the future recording. Our contribution concerns their integration with
structured motion generation rather than a new music representation or forecaster.

\subsection{Structured Motion and Streaming Continuation}
Vector quantization~\cite{vandenOord2017vqvae} supports token-based motion
generation in T2M-GPT~\cite{zhang2023t2mgpt} and MoMask~\cite{guo2024momask},
while latent diffusion~\cite{chen2023mld} provides a continuous alternative.
DisCoRD~\cite{cho2025discord} connects discrete tokens with continuous motion
decoding, and DC-Motion~\cite{wang2026dcmotion} combines discrete structure with
continuous residuals. Our D+C design builds on this representation principle.
Robot-specific reconstruction objectives explicitly supervise designated
structural properties in the discrete branch, while full-motion reconstruction
trains the joint representation through a shared decoder. This encourages
complementary roles rather than assuming that a discrete--continuous split
automatically produces kinematic specialization.

Streaming generation also requires learning to continue from generated context.
Scheduled Sampling~\cite{bengio2015scheduled},
Professor Forcing~\cite{lamb2016professor}, and
Diffusion Forcing~\cite{chen2024diffusionforcing} address training--generation
discrepancies through different treatments of sequence inputs and dynamics.
MotionStream~\cite{jiang2025motionstream} develops causal motion tokenization,
while MotionStreamer~\cite{xiao2025motionstreamer} combines diffusion-based
autoregression with Two-Forward training. \cof{} addresses the transition
induced by our boundary-conditioned D+C generator: sampled structure and detail
enter the committed history together, and their decoded prefix determines the
next reference boundary state. Continuation is trained under this generated
history--state pair, using the reference-update operation employed at inference.

\subsection{Humanoid Motion Generation and Execution}
Motion generation and physical control address complementary aspects of
embodied behavior. PhysDiff~\cite{yuan2023physdiff} incorporates simulator-based
correction into motion generation, while PHC~\cite{luo2023phc} tracks reference
motion with recovery in simulated avatars. ExBody2~\cite{ji2024exbody2},
BeyondMimic~\cite{liao2025beyondmimic}, and SONIC~\cite{luo2026sonic} advance
expressive and versatile humanoid control. Higher-level conditions connect
these capabilities to human intent and perception:
HARMON~\cite{jiang2025harmon} generates humanoid whole-body motion from language,
and RoboPerform~\cite{li2026roboperform} combines motion content with
audio-derived style. Together with DiscoForcing~\cite{ji2026discoforcing},
these systems provide important precedents for condition-driven robot performance.

Reference commitment also relates to receding-horizon action generation in
Diffusion Policy~\cite{chi2023diffusionpolicy} and asynchronous execution in
Real-Time Chunking~\cite{black2025rtc}. The latter freezes actions scheduled for
execution while inpainting the remainder of a new chunk. Such mechanisms govern
how action chunks are updated and delivered; \cof{} instead constructs the
training history and boundary state induced by a joint motion commitment.

Our focus is the kinematically supervised D+C reference generator and its
continuation training. A reference bridge connects its output to a fixed
downstream controller. The generator updates reference-derived context, whereas
the controller closes the measured-state feedback loop. This separation motivates
evaluating musical suitability and motion quality both before and after execution.
